\uuid{qpNX}
\titre{Limite}
\auteur{bodin ; hanani ; mzari}
\organisation{exo7}
\licence{CC BY-NC-SA 4.0}
\datecreate{2020-01-20}
\niveau{L1}
\module{Analyse}
\chapitre{Continuité, limite et étude de fonctions réelles}
\sousChapitre{Limite de fonctions}
\difficulte{2}
\theme{Limites des fonctions réelles, limite}
\exo7qcmid{lille-1.limites.013}

\contenu{
\begin{question}
\qclassification{123.03}
\qtags{motcle=limite}
Soit $f(x)= \frac{\sqrt x}{\sqrt{x+\sqrt{x}}}$. Quelles sont les assertions vraies ?
\begin{answers}

    \good{$\lim_{x\to 0^+} f(x)=0$}

    \bad{$\lim_{x\to 0^+} f(x)=+\infty$}

    \bad{$f$ n'admet pas de limite en $+\infty$.}

    \good{$\lim_{x\to +\infty} f(x)=1$}

\end{answers}
\begin{explanations}
$f(x)= \frac{1}{\sqrt{1+\frac{1}{\sqrt{x}}}}.$
\end{explanations}
\end{question}
}
